%% kapa: exam=1 round=19 year=2008 period=1 form=A subject=appraisal_law subject_ko=부동산관계법규 seq=6 number=46
%% source: Q-Net 감정평가사 기출문제 (cst00309) artlSeq=1001962 file=extracted/1008/1001962/2008년 제19회 감정평가사 1차시험 기출문제/1교시.hwp
%% conversion: kordoc 4.24.51 -> tools/gen1.py (auto); review=proofread-sonnet; flags=none
\begin{question}{46}
국토의 계획 및 이용에 관한 법령의 내용으로 옳은 것은?
\begin{choices}
  \item 도시관리계획결정의 고시가 있는 때에는 국·공유지로서 도시계획시설사업에 필요한 토지를 당해 도시관리계획으로 정하여진 목적외의 목적으로 이를 매각하거나 양도하면 취소사유가 된다.
  \item 도시계획시설사업의 시행자는 사업시행을 위하여 특히 필요하다고 인정되는 때에는 도시계획시설에 인접한 토지·건축물 또는 그 토지에 정착된 물건이나 그 토지·건축물 또는 물건에 관한 소유권외의 권리를 환매부로 수용할 수 있다.
  \item 개발제한구역 안에서의 행위제한 그 밖에 개발제한구역의 관리에 관하여 필요한 사항은 국토의 계획 및 이용에 관한 법률에 의한다.
  \item 국토의 계획 및 이용에 관한 법률 또는 다른 법률에 특별한 규정이 있는 경우를 제외하고는 특별시장·광역시장·시장 또는 군수가 관할구역안의 도시계획시설사업을 시행한다.
  \item 도시계획시설사업의 시행자는 도시계획시설사업의 시행을 위하여 필요한 때에는 등기소 그 밖의 관계 행정기관의 장에게 수수료를 지급하고 필요한 서류의 열람 또는 복사나 그 등본 또는 초본의 교부를 청구할 수 있다.
\end{choices}
\end{question}
\answer{4}
